\documentclass{article}

\usepackage{geometry}
\usepackage{url}
\usepackage{amsmath}
\usepackage{graphicx}

\geometry{left=3cm,right=3cm,top=2.5cm,bottom=2.5cm}

\begin{document}

\begin{titlepage}
    \centering
    \vspace*{6cm} % Adjust the vertical spacing as needed
    {\Huge \textbf{Prelab 0}} \par
    \vspace*{1cm}
    {\Large Runzhao Guo} \par
    \vspace*{1cm}
    {\Large \today} \par
    \vspace*{2cm}
    % Add a logo or image here using \includegraphics
\end{titlepage}

\section{Basics of optical alignment}

\subsection{task 1.(a)}
{The standard for laser classification changed over time, the currently accepted version is one that 
specified by IEC 60825-1 standard}\par
{Class 1 : MPE cannot be exceeded when viewing with naked eyes or through any type of optics.}\par
{Class 1M : Safe in all conditions except through magnifying optics}\par
{Class 2 : Since blink reflex limits the exposure time to laser under 0.25 seconds
\cite{enwiki:1248185910}. It only apply to visible light laser. For continuous wave laser, the power
 is limited to 1mW.}\par
{Class 3R : MPE is exceeded but risk of injury remained low. For visible continuous wave,the power 
is limited to 5mW}\par
{Class 3B : Hazardous to eyes if exposed directly, its diffused reflection, however, is not
For continuous wave with wavelength from 315nm to far infrared, the power is limited to 0.5w. For pulsed 
laser in visible light, the energy is limited to 30mJ.}\par
{Class 4 : All lasers exceeding the Class 3B AEL are Class 4 Lasers.}\par
{Optical Density : The measurement of how much light can protective optics attenuate.}\par
{MPE:the aximum amount of laser can a person take before suffering from any form of injury.}\par

\subsection{1.(b)}
{the document \{laser+MPE...calculation\} provide an equation that is said to be suitable for visible lasers in Example 
1.1, where it mentioned a Table 3, which I can't find. According to the equation}\par
$$
MPE = 1.8t^{\frac{3}{4}}\times 10^{-3}J\cdot cm^{-2}
$$
{radiant exposure $H = 0.636 mJ\cdot cm^{-2}$}\par
{in terms of irradiance, 
\begin{equation}
    \begin{split}
        E =& \frac{H}{t}\\
        =&2.55mW\cdot cm^{-2}
    \end{split}
\end{equation}}\par

\subsection{1.(c)}
{according to the video, cornea, retina and the lens are most susceptible to damage.}\par

\subsection{1.(d)}
{according to the Wikipedia page given in the asignment, visible to near-infrared is the most dangerous to the eye}\par

\subsection{1.(e)}
{choose goggles according to the specific work mode, wave length range, maximum laser power and coresponding minimum 
portection level needed for the given power}\par

\subsection{task 2}
{Fresnel reflection occurs when laser encounter impedance discontinuity, i.e. entering lenses made of fused silica. The amount of 
laser reflected is discribed by reflecance, calculated by using the following equation:}\par
$$
R=(\frac{n_{2}-n_1}{n_2+n_1})^2
$$
{where $n_1$ and $n_2$ are the refractiveindex of the mediums int this task, $n_1\approx 1$ and $n_2\approx 1.458$
therefore}\par
$$R=0.042$$
{this means $4.2\% $ of the light is reflected while $95.8\% $ of the light when in the fused silica lens.}\par
{Now consider the question when a ray of light went through the lens. the discontinuity of impedance occurs twice: both entering 
and exiting the lens. So the loss during the process can be calculated as such:}\par
$$
1-(1-R)^2
$$
{the result is: $9.75\% $ of the light is lost going through the lens.}\par

\subsection{task 3}
{the power of the laser can be calulatd as such:}\par
$$
P=\frac{V_{load}}{R_{load}}R_{DET36A2}
$$
{according to the product page of DET36A2, in the conditions given in the question, $R_{DET36A2}\approx 0.3A/W $ \cite{DET36A2}, 
therefore the power of the laser is:}\par
$$
P \approx 12 \mu W
$$

\section{Gaussian Beams and Telescopes}
\subsection{task 4}
{The ABCD matrix of the objective lens, the gap between lenses and the eyepiece lens are given as follows:}\par
$$
M_{objective}=
\left[ {\begin{array}{cc}
    1 & 0 \\
    -\frac{1}{f_o} & 1 \\
\end{array} } \right]
$$
$$
M_{gap}=
\left[ {\begin{array}{cc}
    1 & f_o + f_e \\
    0 & 1 \\
\end{array} } \right]
$$
$$
M_{eyepiece}=
\left[ {\begin{array}{cc}
    1 & 0 \\
    -\frac{1}{f_e} & 1 \\
\end{array} } \right]
$$
{the ABCD matrix of the telescope can be given by:}\par
\begin{equation}
    \begin{split}
        M_{telescope} &=
        \left[ {\begin{array}{cc}
            1 & 0 \\
            -\frac{1}{f_e} & 1 \\
        \end{array} } \right]
        \left[ {\begin{array}{cc}
            1 & f_o + f_e \\
            0 & 1 \\
        \end{array} } \right]
        \left[ {\begin{array}{cc}
            1 & 0 \\
            -\frac{1}{f_o} & 1 \\
        \end{array} } \right]\\
        &=
        \left[ {\begin{array}{cc}
            -\frac{f_e}{f_o} & f_o + f_e \\
            0 & -\frac{f_o}{f_e} \\
        \end{array} } \right]
    \end{split}
\end{equation}
{The vector of the input light ray that we actually care:
\[
\begin{pmatrix}
w\\
0\\
\end{pmatrix}
\]
and the output vector is
\[
\begin{pmatrix}
-\frac{f_e}{f_o}w\\
0\\
\end{pmatrix}
\]
it is clear that the waist of output ray is $\frac{f_e}{f_o}w$, to answer the question, the ratio of input 
beam waist to output beam wasit is$\frac{f_o}{f_e}$.}

\subsection{task 5}
{I am finding this question a little strange, maybe I'm not understanding it corectly.}\par
{When we are looking into the telescope, our eyes are capturing the light coming out of the eyepiece lens, 
which is parallel, as stated in the question. When we are looking at the world using our naked eyes, our eyes 
are also capturing parallel light rays. We can see images with our naked eyes, so it is natural for our eyes 
to be able to see images from the telescope, because they are taking in input of the same nature.}\par
{Simply put: naked eyes can translate parallel light rays into image, the light rays coming out of telescope are 
parallel, therefore we can see image by looking into a telescope}\par
{As for why naked eyes can translate parallel light rays into image, it is because the optical lens in our eyeballs 
can focus parallel rays of light onto our retina, where light wave of certain wavelength is translated into some sort 
of signal apprehensible by human brain.}\par

\subsection{task 6}
{The magnification M of the telescope is given by the ratio of the output ray angle to the input ray angle. We can easily 
derive from the ABCD matrix obtained in task 4 that the absolute value of the ratio of the output to input angles 
is $\frac{f_o}{f_e}$}\par

\subsection{task 7}
{the modification I'd make is as the picture depicted below:}\par
\begin{figure}[htbp]
    \centering
    \includegraphics[width=0.5\textwidth]{t7.jpg}
    \caption{modification that allow CCD to capture the image}
    \label{t7}
\end{figure}
{the ABCD matrix of the system with the added lens is}\par
\begin{figure}[htbp]
    \centering
    \includegraphics[width=0.5\textwidth]{t7_2.jpg}
    \caption{modification that allow CCD to capture the image}
    \label{t7_2}
\end{figure}
{the magnification changed after my modification}

\bibliographystyle{plain}
\bibliography{references}
\end{document}